\runninghead{The Watch at the Door}
\properlane{watch-at-the-door}{The Watch at the Door: The Hidden Hour and the Command to All}
\label{sec:watch-at-the-door}

The second question comes from the Gospel. Jesus has just said that no one
knows the day or the hour, ``neither the angels in heaven, nor the Son, but
the Father'' (Mk 13:32), and then commands his disciples to stay awake for
an hour they cannot know, and extends the command to everyone: ``what I say to
you, I say to all: Watch''. Two difficulties follow. Why should the hour be
hidden at all? And how can a command to wait for the end be addressed to
\emph{all}, when nearly everyone who hears it will die long before the end
comes? Jerome and John Chrysostom, commenting on the same discourse in
Matthew, answer the first: the hour is hidden for the disciples' good.
Augustine, replying to a bishop's letter, answers the second. Gregory the
Great and Bede say what keeping awake is and who must keep it. The command
to watch therefore holds for the whole time between the two comings: always, because the hour is hidden, and for everyone, because
each one's death brings the end to him.

\subsection{Jerome: you do not know}

Jerome's commentary on Matthew meets the
saying about the hour first as a controversy. The Arians, who held the Son to
be a creature, used it to prove that he was ignorant of something the Father
knew; Jerome answers that the Word, through whom all times were made, the day
of judgment among them, cannot be ignorant of one day in the times he made.
The treasures of wisdom are hidden in Christ (Col 2:3), and after his
resurrection he told the apostles, ``It is not for you to know the times''
(Acts 1:7). He knows, Jerome concludes, but it does not profit the apostles to
know: \latin{ut semper incerti de adventu Judicis, sic quotidie vivant, quasi
die alia judicandi sint}, ``so that, always uncertain of the Judge's coming,
they may live each day as if they were to be judged the next''
(\work{Commentarii in Matthaeum} IV, on Mt 24:36, PL~26, col.~181).

A few verses later Matthew gives the command to watch, and Jerome says that it
shows plainly why the Lord had refused the date: \latin{ut pendulae
exspectationis incerto semper eum credant esse venturum, quem ignorant quando
venturus sit}, ``so that, in the uncertainty of a suspended expectation, they
may always believe that he will come, since they do not know when he will
come''. And he notices the pronoun: \latin{Et non dixit, quia nescimus qua
hora venturus sit Dominus; sed nescitis}, ``He did not say, `we do not know
at what hour the Lord will come', but `you do not know''' (on Mt 24:42--43,
col.~182). The ignorance is the disciples', not the Son's. Bede, commenting
on Mark three centuries later, copied both of Jerome's comments nearly word
for word under Mark 13:32--33 without naming him.

John Chrysostom, preaching on Matthew, makes the same observation in Greek: ``He said not, `I know not,' but, `ye know not'''.
He gives the same reason, and draws from it a description of the life it
asks for: ``For this intent He tells them not, in order that they may watch,
that they may be always ready \ldots\ desiring that they should be anxiously
waiting, and continually in virtuous action.'' Then he widens it to each
person's death: ``if the common sort of men knew when they were to die, they
would surely strive earnestly at that hour. In order therefore that they may
strive, not at that hour only, therefore He tells them not either the common
hour, or the hour of each \ldots\ Wherefore He made the end of each man's
life also uncertain'' (\work{Homilies on Matthew} 77.2--3).

In Jerome's words the disciples are uncertain of the hour and not of the
coming: they ``always believe that he will come''. The Entrance Antiphon
sings a waiting of that kind: none who wait for God (\latin{qui te
exspectant}) will be confounded. Robert Bellarmine, explaining the psalm
verse, says that to wait on the Lord ``is a very common expression in the
Scriptures, and means to expect him in the certain hope of assistance''
(\work{Explanatio in Psalmos} 24, expl.~2--3).

\subsection{Augustine: to all, because the day comes to each}

Chrysostom's turn to the end of each man's life leads to the second
difficulty, and Augustine meets it directly. In a letter to Hesychius, bishop of Salona, who had
asked him about the end of the world, he takes up Mark's last verse and asks
to whom ``what I say to you, I say to all'' was said. Not only to the four on
the Mount of Olives, he answers, but to the elect who belong to the body of
Christ: to those who heard it then, to those after them, to us, and to those
who will come after us until his last coming. The day will not find all of
them alive. The command still binds them all, because \latin{in quo enim
quemque invenerit suus novissimus dies, in hoc eum comprehendet mundi
novissimus dies}: ``in whatever state each one's own last day finds him, in
that state the last day of the world will take him''. Each will be judged on
that day as he was when he died, and therefore every Christian must watch
(\work{Epistula} 199.1.2--3). In the same letter Augustine distinguishes
ignorance of the date, which is no fault, from the bad servant's ``My lord is
long a coming'' (Mt 24:48), which is. Chrysostom, on the second reading, puts
Paul's word to a church impatient at the delay: ``Why make ye much ado
\ldots\ why are ye troubled that Christ is not come? Nay, he is come; and the
Day is henceforth at the doors'' (\work{Homilies on First Corinthians} 2).

Aquinas makes the same connection in commenting on the second reading. Paul
says that Christ will confirm the Corinthians ``unto the end without crime,
in the day of the coming of our Lord'' (1~Cor 1:8). Aquinas reads the end as
the end of each one's life, takes ``without crime'' to mean without mortal
sin, and gives the reason: \latin{qui sine crimine invenitur in die mortis,
sine crimine perveniet ad diem iudicii}, ``whoever is found without crime on
the day of death will come without crime to the day of judgment''
(\work{Super I ad Corinthios} c.~1, lect.~1).

\subsection{Gregory: the knock and the watches}

Gregory preached on Luke's parable of the servants who wait, lamps lit, for
their master to return from a wedding (Lk 12:35--40), a parallel to Mark's
doorkeeper. He reads the master's going to the wedding as the Ascension, and
his return as the judgment. Then he asks what the master's knock is, and
answers: \latin{Venit quippe Dominus cum ad judicium properat, pulsat vero,
cum jam per aegritudinis molestias esse mortem vicinam designat}, ``The Lord
comes when he hastens to judgment; he knocks when, through the troubles of
sickness, he signals that death is near.'' The servant who opens to him at
once is the one who receives him with love; the one who does not want to
open fears to leave the body and to see a Judge he remembers having despised.
``But he who is secure in his hope and his work opens at once to the one who
knocks, for he waits gladly for the Judge'' (\work{Homiliae in Evangelia}
13.3).

Gregory then says what keeping awake is: \latin{Vigilat qui ad aspectum veri
luminis mentis oculos apertos tenet, vigilat qui servat operando quod credit,
vigilat qui a se torporis et negligentiae tenebras repellit}: ``He keeps
watch who holds the eyes of his mind open to the sight of the true light; he
keeps watch who keeps by his works what he believes; he keeps watch who
drives from himself the darkness of sloth and negligence'' (13.3). The
second of these, keeping by works what one believes, is close to what the
Collect asks: the will to go out with just works to meet Christ as he comes.
The Prayer over the Offerings asks that the devotion God grants in this
passing time become for those who offer it the reward of eternal redemption.

Isaiah's people confess the opposite of this wakefulness: ``There is none
that calleth upon thy name: that riseth up, and taketh hold of thee'' (64:6;
Douay 64:7); the Vulgate's verb is \latin{consurgat}, to rise up. In the
psalm's last strophe the people promise the reverse: ``we depart not from
thee \ldots\ and we will call upon thy name''. Augustine explains the words
between, ``Thou wilt quicken us'': ``for dead we were, when to earthly things
we did cleave'', and ``\,`And Thy Name we will invoke;' that is, Thee we will
love'' (\work{Enarrationes in Psalmos} 79.11).

Luke's parable speaks of a master who may come in the second or third watch,
and Gregory reads the watches as the ages of a life: the first is childhood,
the second youth, the third old age. A man who did not keep awake in the
first should keep the second, and one who failed in the second should not
neglect the remedies of the third; \latin{neque sic desperandum est}, ``even
so, there is no cause for despair'' (13.5). His reason for the hidden hour
is Jerome's in other words: \latin{Horam vero ultimam Dominus noster idcirco
voluit nobis esse incognitam, ut semper possit esse suspecta, ut dum illam
praevidere non possumus, ad illam sine intermissione praeparemur}: ``Our Lord
willed that the last hour be unknown to us so that it might always be
suspected, and that, since we cannot foresee it, we may prepare for it
without ceasing'' (13.6). Gregory's three watches are his own allegory on
Luke's verse; Mark's four night watches are a different list. From the hidden
hour Gregory turns to his hearers: since death is certain for all, he says,
do not let the care of earthly things weigh you down, \latin{Nolite ergo attendere
quid habetis, sed quid estis}, ``do not consider what you have, but what you
are'' (13.6). The Prayer after Communion says that by the mysteries God
already teaches those who live among things that pass to love what is
heavenly and hold to what lasts.

\subsection{Bede: the doorkeeper, the servants and every heart}

In Bede's commentary on Mark, what he says on the parable itself is largely
his own. The man gone abroad is Christ, who ascended to the Father after his resurrection and left
the Church in bodily presence, \latin{quam tamen nunquam divinae praesidio
praesentiae destituit}, ``yet never left it without the protection of his
divine presence''. (The identification of the traveller with Christ, and the
sentence explaining that the proper place of the flesh is the earth, Bede
takes from Gregory's homily on the talents.) The master gave his servants
authority over every work, \latin{quia fidelibus suis, concessa sancti
Spiritus gratia, facultatem tribuit bonis operibus serviendi}: ``because, by
granting the grace of the Holy Spirit, he gave his faithful the power to
serve in good works''. He commanded the doorkeeper to watch, ``because he bids
the order of pastors and spiritual rulers to give careful attention to the
Church entrusted to them''.

On the last verse Bede turns from the doorkeeper to everyone: \latin{Non
apostoli tantum et successores eorum, rectores nimirum Ecclesiae, sed et
omnes vigilare praecipimur}, ``Not only the apostles and their successors,
the rulers of the Church, but all of us are commanded to watch.'' Each house
has a door: \latin{Universi januas cordium nostrorum, ne hostis antiquus mala
suggerendo irrumpat, instanter servare jubemur}, ``All of us are bidden to
guard earnestly the doors of our hearts, lest the old enemy break in by
suggesting evil.'' And the reason is personal: \latin{Unusquisque enim pro se
reddet rationem Deo}, ``for each one will give account to God for himself''
(\work{In Marci evangelium expositio} IV, on Mk 13:34--37). He ends with
Gregory's three sentences on keeping watch. In the Preface all who give
thanks call themselves \latin{vigilantes}, keeping watch, and say that they
dare to wait for what was promised.

Augustine and Bede put the weight in different places. Augustine moves the
command from the end of history into each life, so that the watch is kept
against one's own death. Bede keeps it at the Lord's return and divides it by
office: the pastors at the door of the Church, the faithful at their work,
every person at the door of his own heart. The Gospel has room for both: the
master returns to one house, and each servant in it has his own work.

\subsection{Where this interpretation strains}

The Fathers who give the reason for the hidden hour comment on parallels:
Jerome and Chrysostom on Matthew 24, Gregory on Luke 12. Only Augustine and
Bede comment on Mark's verses themselves, and Bede's comment on Mark 13:33 is
Jerome's on Matthew. Mark's verse 32 states the same hidden hour and his
verse 35 the same command, so their reason fits Mark's words as well. The strongest
objection is Augustine's own move. Mark's Gospel speaks of the master's
return at the end of history, and Bede keeps it there; applying the command
to each one's death changes its object. Augustine's reason meets the
objection: he keeps the last day of the world and says that it reaches each
person in the day of his death, since each will be judged then as he is
found at death.

\subsection{The four senses of this interpretation}

\begin{fourSenses}
\item[Literal.] On the Mount of Olives Jesus tells Peter, James, John and
Andrew, and through them all, to stay awake because they do not know the
time; the time is withheld so that they may live each day as if they were to
be judged the next (Jerome), always ready and continually at work
(Chrysostom).
\item[Allegorical.] The man gone abroad is Christ ascended (Gregory on
Matthew's traveller, and Bede on Mark's), who left the Church in the body and
never without his presence (Bede); the authority given to the servants is the grace of the Holy Spirit
for good works, and the doorkeeper the order of pastors (Bede).
\item[Moral.] Each guards the door of his own heart against the old enemy's
suggestions and will give account for himself (Bede); he keeps by works what
he believes, and if he slept in one watch of life he keeps the next, without
despair (Gregory); he asks, in the Alleluia verse, to be shown God's mercy and
given his salvation. Augustine, preaching on the verse the Communion Antiphon
sings, ``our earth shall yield her fruit'', asks that God give increase to
the seed he has sown in his hearers (\work{Enarrationes in Psalmos} 84.15;
§12 in the Oxford translation).
\item[Anagogical.] The master returns at the end of time, and that last day
comes to each person in the day of his death, finding him as he will be
judged (Augustine); whoever is found without grave sin on the day of death
will come without it to the day of judgment (Aquinas), and those who now
keep watch will receive openly what was promised (the Preface); the solemn
blessing, where it is used, asks eternal life for those who await him when he
comes again in majesty.
\end{fourSenses}
